\chapter{ニューロンの型を決めない物質}
\label{sec:othermediators}

第\ref{sec:neurontypes}章では、ニューロンを主な神経伝達物質で分類し、10の型を得た（表\ref{tab:types}）。ここで話を複雑にしておこう。ニューロンが放出するのは主な伝達物質だけではなく、脳では主な伝達物質以外の物質も働いている。それらはネットワークが信号を伝え、蓄える仕方を変える。

アドレスバスとブロードキャストバス（\ref{sec:buses}節）のほかに、3本目のバスがある。\emph{逆向きチャネル}である。いくつかの物質は送り手ではなく\emph{受け手}が放出し、間隙を逆向きに送り手の出力側へ戻り、送り手が次に放出する伝達物質の量を変える。通信ネットワークでいえば、送り手が自分の送信を調整するのに使う受信確認に似ている。ニューロンが結合の重みを変えるときに使うのがまさにこのチャネルである。送り手の出力側はこれを通じて受け手の活動、つまり第\ref{sec:dofa}章の学習則の2番目の因子を知る。

既知の神経伝達物質の完全なリストは長い。100を超える物質があり、その大部分は短いタンパク質の鎖、神経ペプチドである。しかしどれもいくつかのクラスに収まり、クラス内の物質は似た振る舞いをする。主な伝達物質にならない物質を、エンジニアが問うであろう三つの問い、つまり信号がどのバスを通るか、速く効くか、ふつうの符号は何か、とともに表\ref{tab:othermediators}にまとめた。これらはニューロンの型を決めない。主な伝達物質と一緒に放出されるか、ニューロン以外から放出されるか、逆向きに進むかのいずれかだからである。

\begin{table}[t]
\centering
\footnotesize
\setlength{\tabcolsep}{4pt}
\renewcommand{\arraystretch}{1.12}
\begin{tabular}{>{\raggedright}p{2.0cm}>{\raggedright}p{2.6cm}>{\raggedright}p{2.55cm}>{\raggedright}p{1.3cm}>{\centering}p{0.8cm}>{\raggedright\arraybackslash}p{4.2cm}}
\toprule
クラス & 物質 & バス & 速さ & 符号 & 計算での役割 \\
\midrule
アミノ酸 & D-セリン & 局所 & 遅い & AND & グルタミン酸受容体の第二の鍵：「AND」条件 \\
\midrule
モノアミン & 微量アミン & ブロードキャスト & 遅い & $\pm$ & モノアミンチャネルの微調整 \\
\midrule
\multirow{2}{2.0cm}{プリン}
 & ATP & アドレス & 速・遅 & $+$ & 共伝達物質。ニューロンとグリアの間の信号 \\
 & アデノシン & 容積 & 遅い & $-$ & 活動とともに蓄積する：負荷カウンタ\cite{Dunwiddie2001} \\
\midrule
\multirow{8}{2.0cm}{神経ペプチド（100種以上）}
 & エンケファリン、エンドルフィン、ダイノルフィン & 容積 & 遅い & $-$ & 伝達の抑制 \\
 & サブスタンス P & 容積 & 遅い & $+$ & 遅い増強 \\
 & ニューロペプチド Y & 容積 & 遅い & $-$ & 遅い抑制 \\
 & ソマトスタチン & 容積 & 遅い & $-$ & 遅い抑制。GABA と共放出 \\
 & 血管作動性腸管ペプチド（VIP） & 容積 & 遅い & $+$ & 脱抑制：抑制性ニューロンの抑制 \\
 & コレシストキニン & 容積 & 遅い & $\pm$ & 一部の抑制性ニューロンで GABA と共放出 \\
 & オレキシン & ブロードキャスト & 遅い & $+$ & 覚醒モードの安定化 \\
 & オキシトシン、バソプレシン & ブロードキャスト & 遅い & $\pm$ & 行動の大域的な設定 \\
\midrule
\multirow{2}{2.0cm}{気体}
 & 一酸化窒素 NO & 逆向き、容積 & 遅い & $\pm$ & 膜を通り抜ける。重みが変わるときの逆向き信号\cite{Garthwaite2008} \\
 & CO、H$_2$S & 容積 & 遅い & $\pm$ & 役割はほとんど研究されていない \\
\midrule
脂質 & 内因性カンナビノイド（アナンダミド、2-AG） & 逆向き & 遅い & $-$ & 受け手が送り手の放出を減らす：重みのフィードバック\cite{WilsonNicoll2001} \\
\bottomrule
\end{tabular}
\caption{ニューロンの型を決めない物質。\emph{バス}、\emph{速さ}、\emph{符号}は表\ref{tab:types}と同じ。加えて、容積バスは伝達物質が放出場所の周りの細胞間隙に広がること、逆向きは受け手から送り手へ戻ること、局所は放出場所の近くで作用することを表す。「AND」は許可信号である。それだけでは電流を起こさないが、それがないと別の入力が働かない。論理ゲート「AND」の2番目の入力のようなものである。神経ペプチドはほぼ常に主な伝達物質と一緒に、主に高いスパイク頻度のときに放出される\cite{vandenPol2012}。}
\label{tab:othermediators}
\end{table}
